\section[The formula for the n-point function for a KP tau-function]{The formula for the $\boldsymbol{n}$-point function for a KP tau-function}\label{sectionKP}

In this section, we first review the explicit formula (see Theorem~\ref{thm:n-point-KP} below) obtained by the second-named author
for the $n$-point function for an arbitrary KP tau-function in the big cell, and then consider the converse statement that
gives a criterion for a KP tau-function.

Before entering into the details let us introduce some {\it notations}.
Denote by $\mathcal{R}$ a suitable ground ring.
For a formal power series $F({\bf p})\in \mathcal{R}[[{\bf p}]]$,
with ${\bf p}=(p_1,p_2,\dots)$, define the $n$-point function associated with $F({\bf p})$ by
\[
G_n(\xi_1, \dots, \xi_n) = \sum_{k_1,\dots,k_n\ge1} \prod_{i=1}^n \frac{k_i}{\xi_i^{k_i+1}} \cdot
\frac{\pd^n F({\bf p})}{\pd p_{k_1} \cdots \pd p_{k_n}}\biggr|_{{\bf p} = {\bf 0}}, \qquad n\ge0.
\]
By a {\it partition} $\mu=(\mu_1,\mu_2,\dots)$, we mean a sequence of
weakly decreasing non-negative integers with $\mu_k=0$
 for sufficiently large $k$.
 The length $\ell(\mu)$ is the number of the non-zero parts of~$\mu$, the weight $|\mu|:=\mu_1+\mu_2+\cdots$,
 and the multiplicity of $i$ in~$\mu$ is denoted by $m_i(\mu)$.
 The set of all partitions will be denoted by~$\mathcal{P}$, and the set of partitions of weight~$d$ is denoted by $\mathcal{P}_d$.
 The
Schur polynomial $s_{\mu}({\bf p})$ associated to~$\mu\in\mathcal{P}$ is a polynomial in the variables
${\bf p} = (p_1,p_2,\dots)$,
defined by
\beq\label{defSchur1}
s_{\mu}({\bf p}) := \det_{1\leq i,j \leq \ell(\mu)}(h_{\mu_i-i+j}({\bf p})),
\eeq
where $h_j({\bf p})$ are polynomials defined by the generating function
\[
\sum_{j=0}^{\infty}h_j({\bf p})z^j := {\rm e}^{\sum_{k=1}^{\infty} \frac{p_k}{k} z^k}.
\]

One constructive way to describe tau-functions of the KP hierarchy is
to use Sato's Grassmannian. Let
\smash{$z^{1/2}\mathcal{R}\bigl[\bigl[z,z^{-1}\bigr]\bigr]$} be the space of formal series in $z$ of half integral powers.
An~element~$V$ in the big cell $Gr_0$ of Sato's Grassmannian is specified by a sequence of series
\begin{gather*} %\label{eqn:Cell}
\Psi_j(z) = z^{j+1/2} + \sum_{i=0}^\infty A_{i,j} z^{-i-1/2}, \qquad j \geq 0,
\end{gather*}
where the coefficients $A_{i,j}$ are called the {\em affine coordinates} of~$V$. To any element~$V\in Gr_0$, there corresponds
a particular KP tau-function $Z_V$ defined by
\beq\label{fromaffinetotau}
Z_V:= \sum_{\lambda\in\mathcal{P}} c_\lambda s_\lambda({\bf p}),
\eeq
where $T_k=p_k/k$, and for a partition~$\lambda$ written in terms of the Frobenius notation~\cite[Section~I.1, p.~3]{Mac}
$\lambda=(m_1,\dots,m_k|n_1,\dots,n_k)$,
\[
c_\lambda := (-1)^{n_1+\dots+n_k} \det_{1\leq i,j\leq k} \bigl(A_{m_i, n_j}\bigr).
\]
In particular, for $\lambda=\bigl(i+1,1^j\bigr)$ we have $c_\lambda=(-1)^j A_{i,j}$.
So one can easily read off the affine coordinates from the expansion~\eqref{fromaffinetotau}.

The following theorem was obtained in~\cite{Zhou-Emergent}.

\begin{Theorem}[\cite{Zhou-Emergent}]\label{thm:n-point-KP}
The $n$-point function associated with $\log Z_V$ is given by the following formula:
For $n=1$,
\be \label{eqn:G-1}
G_1(\xi) = \sum_{i,j\geq 0} A_{i,j} \xi^{-i-j-2},
\ee
and for $n \ge2 $,
\be \label{eqn:G-n}
G_n(\xi_1, \ldots, \xi_n)
= (-1)^{n-1} \sum_{\text{$n$-cycles}}
\prod_{i=1}^n B\bigl(\xi_{\sigma(i)}, \xi_{\sigma(i+1)}\bigr)
- \frac{\delta_{n,2}}{(\xi_1-\xi_2)^2},
\ee
where $\sigma(n+1)$ is understood as $\sigma(1)$, and
\[
B(\xi_i, \xi_j) = \begin{cases}
%i_{\xi_i, \xi_j}
\dfrac{1}{\xi_i-\xi_j} + A(\xi_i, \xi_j), & i \neq j, \\
A(\xi_i, \xi_i), & i =j, \\
%i_{\xi_j, \xi_i} \frac{1}{\xi_i-\xi_j} + A(\xi_i, \xi_j), & i \neq j,
\end{cases}
\]
and
\[
\begin{split}
A(\xi, \eta)
= & \sum_{i,j \geq 0} A_{i,j} \xi^{-j-1} \eta^{-i-1} .
\end{split}
\]
\end{Theorem}

Theorem~\ref{thm:n-point-KP} enables one to easily compute the $n$-point function once
we have found the affine coordinates $A_{i,j}$.
For example, there is only one $2$-cycle,
so the two-point function is given by
\ben
G_2(\xi_1, \xi_2)
& = & \frac{A(\xi_1, \xi_2) - A(\xi_2,\xi_1)}{\xi_1-\xi_2}
-A(\xi_1, \xi_2)\cdot A(\xi_2,\xi_1).
\een


Let us consider the converse of the above Theorem~\ref{thm:n-point-KP}. Recall the simple fact, which can be verified directly using~\eqref{bilinearidH}, that
multiplying a KP tau-function $\tau({\bf T}({\bf p}))$
by the exponential of an arbitrary affine linear function of~${\bf p}$
\beq\label{linear1004}
\tau({\bf T}({\bf p})) \cdot {\rm e}^{C_0+\sum_{k\ge1} C_k p_k}, \qquad C_k\in\mathcal{R},\quad k\ge0,
\eeq
produces again a KP tau-function. We recall here that $T_k=p_k/k$. The renormalized variables~$p_k$ are also sometimes called KP times.

\begin{Corollary}\label{conversthmA}
Let $F$ be a formal power series in~${\bf p}$ and $G_n(\xi_1,\dots,\xi_n)$ its $n$-point function.
If~there are $(A_{i,j})_{i,j\geq 0}$ such that for all $n\ge2$
$G_n$ are given by~\eqref{eqn:G-n},
then $Z= {\rm e}^F$ is a tau-function of the KP hierarchy.
\end{Corollary}
\begin{proof}
By Theorem~\ref{thm:n-point-KP},
the tau-function corresponds to the point in the big cell with $A_{i,j}$ being the affine coordinates
can only differ from $Z$ by multiplying by the exponential of some affine-linear function.
\end{proof}

\begin{Remark}
From the above proof, we see that Corollary~\ref{conversthmA} directly follows from Theorem~\ref{thm:n-point-KP}. In~\cite{ABDBKS},
Theorem~\ref{thm:n-point-KP} together with Corollary~\ref{conversthmA} is regarded as {\it Zhou's theorem}.
\end{Remark}

The next proposition describes how the affine coordinates change under~\eqref{linear1004}.

\begin{Proposition}\label{VVprimelinear}
Let $V$, $V'$ be two elements in the big cell having affine coordinates~$A_{i,j}$,~$A'_{i,j}$, respectively,
and let $Z_V$, $Z_{V'}$ be the KP tau-functions associated to~$V$, $V'$. If there exist $C_1,C_2,\dots$ such that
$
Z_{V'} = {\rm e}^{\sum_{k\ge1} C_k p_k} Z_V$,
then we have the following identity:
\begin{gather*}%\label{AAprime}
\sum_{i,j \geq 0} A'_{i,j} \xi^{-j-1} \eta^{-i-1} + \frac1{\xi-\eta} = \biggl(\sum_{i,j \geq 0} A_{i,j} \xi^{-j-1} \eta^{-i-1} + \frac1{\xi-\eta}\biggr)
\frac{{\rm e}^{-\sum_{\ell\ge1} C_\ell \xi^{-\ell}}}{{\rm e}^{-\sum_{\ell\ge1} C_\ell \eta^{-\ell}}}.
\end{gather*}
\end{Proposition}
\begin{proof}
By direct verifications using~\eqref{eqn:G-1}.
\end{proof}

\section{From the Toda lattice hierarchy to the KP hierarchy}\label{sectionToda}
In this section, we show that an arbitrary tau-function for the Toda lattice hierarchy is a KP tau-function.

Let $(v(x,{\bf m};\e), w(x,{\bf m};\e))$ be an arbitrary solution in $W[[{\bf m}]]^2$
to the Toda lattice hierarchy~\eqref{TLH}, and $\tau(x,{\bf m};\e)$ the tau-function~\cite{CDZ, DY, DZ-toda} of this solution.
Here, $W$ is a certain ring of functions of~$x$ closed under shifting~$x$
by $\pm \e$, and $\tau(x,{\bf m};\e)$ lives in \smash{$\widetilde W[[{\bf m}]]$} with \smash{$\widetilde W$} being some extension of~$W$.
The solution $(v(x,{\bf m};\e), w(x,{\bf m};\e))$ is in one-to-one correspondence
with its initial value
\[
v(x,{\bf 0};\e)=:f(x,\e), \qquad w(x,{\bf 0}; \e)=:g(x,\e).
\]
Denote by $L_{\rm ini}$ the initial Lax operator
\[
L_{\rm ini}=\Lambda+f(x,\e) + g(x,\e)\Lambda^{-1},
\]
and let
\[
s(x,\e) = - \bigl(1-\Lambda^{-1}\bigr)^{-1} (\log g(x,\e)),
\]
which lives in the extended ring~\smash{$\widetilde{W}$}.
Recall from~\cite{Y} that two elements
\[
 \psi_1(\lambda,x;\e) = \bigl(1+{\rm O}\bigl(\lambda^{-1}\bigr)\bigr)\lambda^{x/\e}, \qquad
\psi_2(\lambda,x;\e)=\bigl(1+{\rm O}\bigl(\lambda^{-1}\bigr)\bigr){\rm e}^{-s(x,\e)}\lambda^{-x/\e}
\]
are called forming {\it a pair of wave functions} of $L_{\rm ini}$ if
\begin{align*}%\label{pairwave1}
& L_{\rm ini} (\psi_1) = \lambda \psi_1, \qquad L_{\rm ini} (\psi_2) = \lambda \psi_2,\\
%\label{pairwave2}
& d(\lambda,x;\e):= \psi_1 \, \Lambda^{-1}(\psi_2) - \psi_2 \, \Lambda^{-1}(\psi_1) = \lambda {\rm e}^{-s(x-\e,\e)}.
\end{align*}
The following formula is proved in~\cite{Y}:
\begin{align}
& \e^n \sum_{i_1,\dots,i_n\ge0}
\frac{(i_1+1)! \cdots (i_n+1)!} {\lambda_1^{i_1+2}\cdots \lambda_n^{i_n+2}} \,
\frac{\p^n \log \tau(x,{\bf m};\e)}{\p m_{i_1} \cdots \p m_{i_n}}\bigg|_{{\bf m}={\bf 0}}\nn\\
&\qquad{}= (-1)^{n-1} \frac{{\rm e}^{n s(x-\e,\e)}}{\prod_{j=1}^n \lambda_j} \sum_{\text{$n$-cycles}}
\prod_{j=1}^n D\bigl(\lambda_{\sigma(j)},\lambda_{\sigma(j+1)}; x;\e\bigr) - \frac{\delta_{n,2}}{(\lambda-\mu)^2},\label{Y421}
\end{align}
where $n\ge2$, and
\beq\label{dpsiabtoda}
D(\lambda,\mu; x;\e) := \frac{\psi_1(\lambda,x;\e) \psi_2(\mu,x-\e;\e) - \psi_1(\lambda,x-\e;\e) \psi_2(\mu,x;\e)}{\lambda-\mu}.
\eeq
\begin{Remark}
The function $D(\lambda,\mu; x;\e)$ looks similar to the
Christoffel--Darboux kernel in~matrix models~\cite{Deift, Mehta}.
We hope to study their relations in a future work.
\end{Remark}

It was observed in~\cite[(99)]{Y} that
\smash{$\frac{{\rm e}^{s(x-\e,\e)}}{\mu} D(\lambda,\mu; x;\e) \frac{\mu^{x/\e}}{\lambda^{x/\e}} - \frac1{\lambda-\mu}$}
is a power series of~$\lambda^{-1}$,~$\mu^{-1}$.
By a more careful analysis we can show that
\beq\label{eqn:hat-A-26}
\frac{{\rm e}^{s(x-\e,\e)}}{\mu} D(\lambda,\mu; x;\e) \frac{\mu^{\frac{x}\e}}{\lambda^{\frac{x}\e}} =
\frac1{\lambda-\mu} + \sum_{i,j\ge0} \frac{A_{i,j}(x,\e)}{\lambda^{j+1} \mu^{i+1}} =: B(\lambda,\mu;x;\e)
\eeq
for some coefficients $A_{i,j}(x,\e)$.
%We note that the expansion given in~\cite{Y} is slightly weaker than~\eqref{eqn:hat-A-26}.
In the next section, we will give a complete proof of a generalized version of~\eqref{eqn:hat-A-26}.

Using~\eqref{eqn:hat-A-26} and~\eqref{Y421} we find
\begin{align}
&\e^n \sum_{i_1,\dots,i_n\ge0}
\frac{(i_1+1)! \cdots (i_n+1)!} {\lambda_1^{i_1+2}\cdots \lambda_n^{i_n+2}} \,
\frac{\p^n \log \tau(x,{\bf m};\e)}{\p m_{i_1} \cdots \p m_{i_n}}\bigg|_{{\bf m}={\bf 0}}\nn \\
&\qquad= (-1)^{n-1} \sum_{\text{$n$-cycles}}
\prod_{j=1}^n B\bigl(\lambda_{\sigma(j)},\lambda_{\sigma(j+1)};\e\bigr) - \frac{\delta_{n,2}}{(\lambda-\mu)^2}, \qquad n\ge2.\label{Y926}
\end{align}
(See also~\cite[Corollary~1]{Y}.)

We are ready to prove Theorem~\ref{thm:toda-kp}.

\begin{proof}[Proof of Theorem~\ref{thm:toda-kp}]
By using \eqref{eqn:hat-A-26}, \eqref{Y926} and~\eqref{tT30}, we know that the $n$-point function for~${n\ge2}$
associated to $\log\tau(x,{\bf m};\e)$ has the form~\eqref{eqn:G-n}.
The theorem is then proved by applying Corollary~\ref{conversthmA}.
\end{proof}

If we write
\[
\tau(x,{\bf m};\e) = \tau_{\rm corr}(x,\e) \tau_1 (x,{\bf m};\e), \qquad \tau_1(x, {\bf m}={\bf 0};\e)\equiv1,
\]
then by Theorem~\ref{thm:toda-kp} we know that $\tau_1(x, {\bf m}; \e)$ is a KP tau-function in the big cell.
The factor~$\tau_{\rm corr}(x, \e)$ can~\cite{CDZ, Y} be determined by
\beq\label{correctionfactor}
\frac{\tau_{\rm corr}(x+\e,\e)\tau_{\rm corr}(x-\e,\e)}{\tau_{\rm corr}(x,\e)^2} = g(x,\e) = {\rm e}^{s(x-\e,\e)-s(x,\e)}.
\eeq
This factor can be identified with the {\it correction factor} in~\cite{DY, Y, YZ}.

Let $\tilde{\tau}_1(x, {\bf m};\e)$ be the KP tau-function
associated to the point in Sato's Grassmannian having the affine coordinates
$A_{i,j}(x,\e)$, with $A_{i,j}(x,\e)$ given in~\eqref{eqn:hat-A-26}. Then there exist
$a_i(x,\e)$, $i\ge0$, such that
\[
\tau_1(x, {\bf m};\e) = {\rm e}^{\sum_{i\ge0} a_i(x,\e) m_i} \tilde{\tau}_1(x, {\bf m};\e).
\]
By further applying Proposition~\ref{VVprimelinear}, we get the affine coordinates for $\tau_1(x, {\bf m};\e)$.

\section{Extended Toda flows in the KP hierarchy}\label{sectionTodaext}
The goal of this section is to prove Theorem~\ref{thm:toda-kp-ext}.

Before entering into the main construction for this section, we briefly recall here the definition of $\log L$.
It is shown in~\cite{CDZ} that there exist dressing operators
\begin{align*}%\label{CoePk1005new}
P= \sum_{k\ge0} P_k \Lambda^{-k}, \qquad P_0= 1,\qquad
Q= \sum_{k\ge0} Q_k \Lambda^k, %\label{CoeQk1005new}
\end{align*}
such that
\[
L=P\circ\Lambda\circ P^{-1}=Q\circ\Lambda^{-1}\circ Q^{-1}.
\]
Here, $P_k$, $Q_k$ belong to a certain extension of the differential polynomial ring.
The logarithm of the Lax operator is then defined by~\cite{CDZ}
\[
\log L= \frac12 P\circ \e\p_x \circ P^{-1} - \frac12 Q\circ \e\p_x \circ Q^{-1}.
\]

As in the introduction, for an arbitrary solution $(v({\bf b}+x{\bf 1}, {\bf m};\e), w({\bf b}+x{\bf 1}, {\bf m};\e))$ to the ETH,
let $\tau({\bf b}+x{\bf 1}, {\bf m};\e)$ be the tau-function (again in the sense of~\cite{CDZ, DZ-toda})
 of the solution. Let us also denote
 \beq\label{defsigma}
 \sigma(x, {\bf b}, {\bf m};\e) = \log\frac{\tau({\bf b}+x{\bf 1}, {\bf m};\e)}{\tau({\bf b}+(x+\e){\bf 1}, {\bf m};\e)}.
 \eeq

\subsection{The matrix-resolvent method}
The matrix-resolvent method for studying tau-functions for the Toda lattice hierarchy was developed in~\cite{DY} (see also~\cite{Y}).
By the locality nature of this method, the same formulation as~for the Toda lattice hierarchy
applies to the ETH.

Denote by
\[
\mathcal{A}=\ZZ[v_0, w_0, v_{\pm1}, w_{\pm 1}, v_{\pm 2}, w_{\pm 2}, \dots]
\]
the polynomial ring.
Recall that the {\it basic matrix resolvent} $R(\lambda)$ is defined~\cite{DY} as
the unique element in ${\rm Mat}\bigl(2,\mathcal{A}\bigl[\bigl[\lambda^{-1}\bigr]\bigr]\bigr)$ satisfying
\begin{align*}
& \Lambda(R(\lambda)) U(\lambda) - U(\lambda) R(\lambda) =0,\\
& \tr \, R(\lambda) =1, \qquad \det R(\lambda)=0,\qquad
 R(\lambda) - \begin{pmatrix} 1 & 0 \\ 0 & 0 \\ \end{pmatrix} \in {\rm Mat}\bigl(2,\mathcal{A}\bigl[\bigl[\lambda^{-1}\bigr]\bigr]\lambda^{-1}\bigr), %\label{entrymr421}
\end{align*}
where
\[
U(\lambda) := \begin{pmatrix} v_0-\lambda & w_0 \\ -1 & 0 \\ \end{pmatrix}.
\]
Define \smash{$S_i=\frac1{(i+1)!}{\rm Coef} \bigl(\Lambda(R(\lambda)_{21}), \lambda^{-i-2}\bigr)\in\mathcal{A}$}, $i\geq0$, and
a sequence of elements $\Omega_{i,j} \in \mathcal{A}$ by
\[
\e^2\sum_{i,j\geq0} \frac{(i+1)!(j+1)!}{\lambda^{i+2} \mu^{j+2}}\Omega_{i,j} = \frac{\tr \, R(\lambda)R(\mu)}{(\lambda-\mu)^2} - \frac{1}{(\lambda-\mu)^2}.
\]

According to~\cite{DY}, the above-defined $(\Omega_{i,j}, S_i)$ gives rise to part of
the canonical tau-structure for the ETH in~\cite{CDZ, DZ-toda}.
Indeed, it is shown in~\cite{DY} that $(\Omega_{i,j}, S_i)$ is associated to the
tau-symmetric hamiltonian densities $h_{\alpha,p}$, $\alpha=1,2$, $p\ge-1$,~\cite{CDZ, DZ-toda} for the Toda lattice hierarchy.
More precisely, $S_i=h_{2,i-1}$ and the following identities hold
for an arbitrary solution $(v({\bf b}+x{\bf 1},{\bf m};\e), w({\bf b}+x{\bf 1},{\bf m};\e))$ to the ETH:
\begin{align}
& \e^2\frac{\p^2 \log\tau({\bf b}+x{\bf 1},{\bf m};\e)}{\p m_i \p m_j}=\Omega_{i,j}({\bf b}+x{\bf 1},{\bf m};\e),\nn \\%\label{deftau1}
& \e \frac{\p}{\p m_i} \biggl(\log \frac{\tau({\bf b}+(x+\e){\bf 1},{\bf m};\e)}{\tau({\bf b}+x{\bf 1},{\bf m};\e)}\biggr) =
S_i({\bf b}+x{\bf 1},{\bf m};\e), \label{deftau2} \\
& \frac{\tau({\bf b}+(x+\e){\bf 1},{\bf m};\e)\tau({\bf b}+(x-\e){\bf 1},{\bf m};\e)}{\tau({\bf b}+x{\bf 1},{\bf m};\e)^2} = w({\bf b}+x{\bf 1},{\bf m};\e). \label{deftau3}
\end{align}
Here, $\Omega_{i,j}({\bf b}+x{\bf 1},{\bf m};\e)$ and
$S_i({\bf b}+x{\bf 1},{\bf m};\e)$ are obtained by replacing $v_k$, $w_k$ by $v({\bf b}+x{\bf 1}+k\e,{\bf m};\e)$ and
$w({\bf b}+x{\bf 1}+k\e,{\bf m};\e)$, respectively.
According to~\cite{DY}, for any $n\ge2$,
\begin{align}
&\e^n \sum_{i_1,\dots,i_n\ge0}
\prod_{j=1}^n \frac{(i_j+1)!}{\lambda_j^{i_j+2}} \frac{\p^n \log \tau({\bf b}+x{\bf 1},{\bf m}; \e)}{\p m_{i_1} \cdots \p m_{i_n}}\nn\\
&\qquad{}= - \sum_{\text{$n$-cycles}}
\frac{\tr \, \prod_{j=1}^n R(\lambda_{\sigma(j)})}{\prod_{j=1}^n \bigl(\lambda_{\sigma(j)}-\lambda_{\sigma(j+1)}\bigr)} - \frac{\delta_{n,2}}{(\lambda-\mu)^2}.\label{multiderivativesresolvent}
\end{align}

Note that by~\eqref{defsigma} and \eqref{deftau2}, \eqref{deftau3}, we have
\begin{align}
& - \e\frac{\p \sigma(x, {\bf b},{\bf m};\e)}{\p m_i} = S_i({\bf b}+x{\bf 1},{\bf m};\e), \label{property1sigma} \\
& {\rm e}^{\sigma(x-\e, {\bf b},{\bf m};\e)-\sigma(x, {\bf b},{\bf m};\e)} = w({\bf b}+x{\bf 1},{\bf m};\e). \label{propertysigma}
\end{align}
It is also clear from~\eqref{correctionfactor} that we can choose
$\sigma(x, b_0=0, b_1=0, \dots, {\bf m}={\bf 0};\e)=s(x,\e)$.

\subsection{Wave functions for ETH and the KP hierarchy}
It is shown in~\cite{Carlet} (see also~\cite{CL, Milanov}) that there exist
dressing operators $P$, $Q$ of the forms
\begin{align}
&P= \sum_{k\ge0} P_k(x, {\bf b}, {\bf m};\e) \Lambda^{-k}, \qquad P_0(x, {\bf b}, {\bf m};\e)\equiv 1, \label{CoePk1005}\\
&Q= \sum_{k\ge0} Q_k(x, {\bf b}, {\bf m};\e) \Lambda^k, \label{CoeQk1005}
\end{align}
such that
\begin{align*}
& L=P\circ\Lambda\circ P^{-1}=Q\circ\Lambda^{-1}\circ Q^{-1}, \\
& \e\frac{\p P}{\p m_i} = - \frac1{(i+1)!} \bigl(L^{i+1}\bigr)_- \circ P, \\
& \e\frac{\p P}{\p b_i} = - \frac{2}{i!}\bigl(L^{i}(\log L-c_i)\bigr)_- \circ P,\\
& \e\frac{\p Q}{\p m_i} = \frac1{(i+1)!}\bigl(L^{i+1}\bigr)_+ \circ Q, \\
& \e\frac{\p Q}{\p b_i} = \frac{2}{i!}\bigl(L^{i}(\log L-c_i)\bigr)_+ \circ Q.
\end{align*}
Define
\begin{align*}
& \psi_1(\lambda; x, {\bf b}, {\bf m};\e) :=
P\Bigl(\lambda^{\frac{x}\e} {\rm e}^{\sum_{i\ge0} \frac{2}{i!} \frac{b_i}\e \lambda^i (\log \lambda-c_i)+\sum_{i\ge0} \frac{m_i}{(i+1)!\e} \lambda^{i+1}}\Bigr) , \\
& \psi_2(\lambda; x, {\bf b}, {\bf m};\e) := Q\Bigl(\lambda^{-\frac{x}\e}{\rm e}^{-\sum_{i\ge0} \frac{2}{i!} \frac{b_i}\e \lambda^i (\log \lambda-c_i)-\sum_{i\ge0} \frac{m_i}{(i+1)!\e} \lambda^{i+1}}\Bigr).
\end{align*}
Then $\psi_a=\psi_a(\lambda; x, {\bf b}, {\bf m};\e)$, $a=1,2$, satisfy the following equations:
\begin{align}
& L\psi_a = \lambda \psi_a, \qquad a=1,2, \label{Laxsp1005}\\
& \e \frac{\p \psi_1}{\p m_i} = \frac1{(i+1)!} \bigl(L^{i+1}\bigr)_+ \psi_1,\nn \\
& \e \frac{\p \psi_1}{\p b_i} = \frac{2}{i!} L^i(\log L - c_i)_+ \psi_1, \label{Laxa1bi} \\
& \e \frac{\p \psi_2}{\p m_i} = -\frac1{(i+1)!} \bigl(L^{i+1}\bigr)_- \psi_2,\nn \\
& \e \frac{\p \psi_2}{\p b_i} = -\frac{2}{i!} L^i(\log L - c_i)_- \psi_2. \label{Laxa2bi}
\end{align}
Moreover, $\psi_1$, $\psi_2$ have the form
\begin{align}
&\psi_1 = \phi_1(\lambda; x, {\bf b}, {\bf m};\e) \lambda^{\frac{x}\e}
 {\rm e}^{\sum_{i\ge0} \frac{2}{i!} \frac{b_i}{\e} \lambda^i (\log \lambda-c_i)+\sum_{i\ge0} \frac{m_i}{(i+1)! \e} \lambda^{i+1}} , \label{psiphiform1} \\
&\psi_2 = \phi_2(\lambda; x, {\bf b}, {\bf m};\e) {\rm e}^{-\sigma(x,{\bf b},{\bf m};\e)} \lambda^{-\frac{x}\e}
 {\rm e}^{-\sum_{i\ge0} \frac{2}{i!} \frac{b_i}{\e} \lambda^i (\log \lambda-c_i)-\sum_{i\ge0} \frac{m_i}{(i+1)! \e} \lambda^{i+1}} , \label{psiphiform2}
\end{align}
where
\begin{align}
& \phi_1(\lambda; x, {\bf b}, {\bf m};\e) =\sum_{k\ge0} P_k(x, {\bf b}, {\bf m};\e) \lambda^{-k}, \label{phiexpansion1}\\
& \phi_2(\lambda; x, {\bf b}, {\bf m};\e) =\sum_{k\ge0} B_k(x, {\bf b}, {\bf m};\e) \lambda^{-k}. \label{phiexpansion2}
\end{align}
We recall that $P_k(x, {\bf b}, {\bf m};\e)$ are given in~\eqref{CoePk1005},
and we also note using \eqref{CoeQk1005} that $B_k(x, {\bf b}, {\bf m};\e)=Q_k(x, {\bf b},{\bf m};\e) {\rm e}^{\sigma(x, {\bf b},{\bf m};\e)}$, $k\geq 0$,
and $B_0(x, {\bf b}, {\bf m};\e)\equiv 1$.
We call $\psi_1$ {\it the wave function of type~A} and $\psi_2$ {\it the wave function of type~B},
{\it associated to the solution $(v(x, {\bf b}, {\bf m};\e), w(x, {\bf b}, {\bf m};\e))$}.

Denote
\begin{align*}
&d(\lambda; x, {\bf b}, {\bf m};\e) \\
&\qquad{}:= \psi_1(\lambda; x, {\bf b}, {\bf m};\e) \psi_2(\lambda; x-\e, {\bf b}, {\bf m};\e)
- \psi_1(\lambda; x-\e, {\bf b}, {\bf m};\e) \psi_2(\lambda; x, {\bf b}, {\bf m};\e).
\end{align*}
The following lemma is important.
\begin{Lemma}\label{keylemma}
For any $i\ge0$, we have the identities
\begin{align}
& \frac{\p \bigl({\rm e}^{\Lambda^{-1}(\sigma(x,{\bf b},{\bf m};\e))}d(\lambda; x, {\bf b}, {\bf m};\e)\bigr)}{\p b_i} =0 , \label{id1} \\
& \frac{\p \bigl({\rm e}^{\Lambda^{-1}(\sigma(x,{\bf b},{\bf m};\e))}d(\lambda; x, {\bf b}, {\bf m};\e)\bigr)}{\p m_i} =0. \label{id2}
\end{align}
\end{Lemma}
\begin{proof}
We first introduce some notations,
\begin{align*}
& R_j := {\rm res} L^j, \qquad N_j={\rm Coef}\bigl(L^j, \Lambda^{-1}\bigr), \\
& r_j:= {\rm res} \log L \circ L^j, \qquad n_j= {\rm Coef}\bigl(\log L \circ L^j, \Lambda^{-1}\bigr).
\end{align*}
({\it Warning: avoid from confusing the notation $R_j$ with the notation for the matrix resolvent.})
For example~ (see \cite{Carlet}),
\begin{align*}
& R_0=1, \qquad R_1=v, \qquad N_0=0, \qquad N_1=w, \\
& r_0= \frac12 \Lambda \circ \frac{\e\p}{\Lambda-1} (v), \qquad n_0 = \frac12 \frac{\e\p}{\Lambda-1} (\log w) .
\end{align*}
Here we omitted the arguments in $R_j$, $r_j$, $N_j$, $n_j$, $v$, $w$. Below, when
no ambiguity will occur, we often do this type of omission.

As it was given in~\cite{CDZ}, the tau-symmetric hamiltonian densities $h_{\alpha, p}$ for the ETH
are related to $R_j$, $r_j$ by
\begin{align}
& h_{2, p} = \frac1{(p+2)!} R_{p+2}, \label{CDZtausymm2}\\
& h_{1,p} = \frac2{(p+1)!} r_{p+1} - \frac{2c_{p+1}}{(p+1)!} R_{p+1}, \qquad p\ge-1.\nn
\end{align}
We also recall from~\cite{CDZ} the following formula:
\[
h_{\alpha,i-1} = \e(\Lambda-1) \frac{\p \log \tau}{\p t^{\alpha,i}} = - \e\frac{\p \sigma}{\p t^{\alpha,i}}, \qquad i\ge0.
\]

Let us now prove~\eqref{id2}. (Although \eqref{id2} was already proved in~\cite{Y},
our proof here will be slightly different from~\cite{Y} and can be generalized to proving~\eqref{id1}.)
We have
\beq\label{derivativemi1005}
{\rm e}^{-\Lambda^{-1}(\sigma)}\frac{\p \bigl({\rm e}^{\Lambda^{-1}(\sigma)}d\bigr)}{\p m_i}
= \frac{\p \Lambda^{-1}(\sigma)}{\p m_i} d + \frac{\p d}{\p m_i}, \qquad i\ge0.
\eeq
Denote
\[
B_i:= -\bigl(L^{i+1}\bigr)_-, \qquad i\ge-1.
\]
Then for all $i\ge0$, we have
\begin{align*}
B_i & = -\bigl(L^i\circ L\bigr)_- = - \bigl(\bigl(L^i\bigr)_+ + \bigl(L^i\bigr)_-\bigr) \circ L_-
 = - (R_i L)_- - \bigl(\bigl(L^i\bigr)_- \circ L\bigr)_- \nn\\
& = B_{i-1} \circ L + N_i - R_i w \Lambda^{-1} = \cdots
 = \sum_{j=0}^i \bigl(N_j-R_j w\Lambda^{-1}\bigr) \circ L^{i-j}. \nn
\end{align*}
(The resulting equality is also valid for $i=-1$ as $B_{-1}=0$.)
Thus,
\beq\label{psi2mi1005}
(i+1)! \e \frac{\p \psi_2}{\p m_i} = B_i (\psi_2) = \sum_{j=0}^i \lambda^{i-j} \bigl(N_j-R_j w\Lambda^{-1}\bigr) (\psi_2),
\eeq
and
\beq\label{psi1mi1005}
(i+1)! \e \frac{\p \psi_1}{\p m_i} %= (L^{i+1} + B_i) (\psi_1)
=\lambda^{i+1} \psi_1
+ \sum_{j=0}^i \lambda^{i-j} \bigl(N_j-R_j w\Lambda^{-1}\bigr) (\psi_1).
\eeq
By substituting \eqref{psi1mi1005} and \eqref{psi2mi1005} in~\eqref{derivativemi1005}
and by a lengthy but straightforward calculation using also~\eqref{Laxsp1005} and \eqref{property1sigma}, we obtain
\begin{align*}
{\rm e}^{-\Lambda^{-1}(\sigma)} \e\frac{\p \bigl({\rm e}^{\Lambda^{-1}(\sigma)}d\bigr)}{\p m_i}
={}& {-}{} \Lambda^{-1}(S_{i}) d
+ \frac{\Lambda^{-1}(R_{i+1})}{(i+1)!} d \nn\\
&{}{+}{} \frac{d}{(i+1)!}\sum_{j=0}^i \lambda^{i-j} \bigl(N_j+\Lambda^{-1}(N_j)+\Lambda^{-1}(v R_j-R_{j+1})\bigr) = 0. \nn
\end{align*}
Note that in the last equality we used \eqref{deftau2}, \eqref{CDZtausymm2}
and the relation $L^{k+1}=L\circ L^k=L^k \circ L$ (which implies the vanishing of each
summand in the above $\sum_j$).

We proceed to prove~\eqref{id1}. Like in~\cite{Y}, we note that
\begin{align*}
&\Lambda \bigl({\rm e}^{\Lambda^{-1}(\sigma(x,{\bf b},{\bf m};\e))}d(\lambda; x,{\bf b}, {\bf m};\e)\bigr)
 = {\rm e}^{\sigma} (\Lambda(\psi_1) \psi_2
- \psi_1 \Lambda( \psi_2) ) \nn\\
&\qquad{} = {\rm e}^{\sigma} \bigl\{ \bigl((\lambda-v) \psi_1 - w \Lambda^{-1}(\psi_1)\bigr) \psi_2 - \psi_1 \bigl((\lambda-v) \psi_2 - w \Lambda^{-1}(\psi_2)\bigr) \bigr\} \nn\\
&\qquad{} = {\rm e}^{\Lambda^{-1}(\sigma(x, {\bf b},{\bf m};\e))}d(\lambda; x,{\bf b}, {\bf m};\e) \nn
\end{align*}
(the last equality used~\eqref{propertysigma}).
It follows that
\[
\e\p_x \bigl({\rm e}^{\Lambda^{-1}(\sigma(x,{\bf b},{\bf m};\e))}d(\lambda; x, {\bf b}, {\bf m};\e)\bigr)
= \frac{\e\p_x}{\Lambda-1} \circ (\Lambda-1) \bigl({\rm e}^{\Lambda^{-1}(\sigma(x, {\bf b},{\bf m};\e))}
d(\lambda; x, {\bf b}, {\bf m};\e)\bigr) = 0.
\]
Thus,
\begin{align}
0 & = {\rm e}^{-\Lambda^{-1}(\sigma(x,{\bf b},{\bf m};\e))} \p_x \bigl({\rm e}^{\Lambda^{-1}(\sigma(x,{\bf b},{\bf m};\e))}d(\lambda; x, {\bf b}, {\bf m};\e)\bigr) \nn\\
& = - \Lambda^{-1}(h_{1,-1}) + \p_x(d(\lambda; x, {\bf b}, {\bf m};\e)) = - 2\Lambda^{-1}(r_0)d + \p_x(d(\lambda; x, {\bf b}, {\bf m};\e))\nn.
\end{align}
Denote
\[
Q_i = -2 \bigl(L^i(\log L-c_i)\bigr)_- = -2 \bigl(L^i \log L\bigr)_0-2c_i B_{i-1}, \qquad i\ge0.
\]
In the way similar to the derivation for $B_i$, we find
\[
Q_i= \sum_{j=0}^{i-1} \bigl(2m_j-2r_jw\Lambda^{-1}\bigr) \circ L^{i-j} - 2 (\log L)_- \circ L^i - 2c_i B_{i-1}.
\]
Therefore,
\begin{align*}%\label{psi2bi1005}
i! \frac{\p \psi_2}{\p b_i} = Q_i (\psi_2)={}&{}
 \sum_{j=0}^{i-1} \lambda^{i-j} \bigl(2m_j-2r_jw\Lambda^{-1}\bigr) (\psi_2)\\
&{}{-}\, 2 (\log L)_- (\psi_2)
-2c_i\sum_{j=0}^{i-1} \lambda^{i-1-j} \bigl(N_j-R_j w\Lambda^{-1}\bigr) (\psi_2),\nn
\end{align*}
and
\begin{align*}%\label{psi1bi1005}
i! \frac{\p \psi_1}{\p b_i} %= (L^{i+1} + B_i) (\psi_1)
={}&{}2 \lambda^i(\log\lambda-c_i) \psi_1+ \sum_{j=0}^{i-1} \lambda^{i-j} \bigl(2m_j-2r_jw\Lambda^{-1}\bigr) (\psi_1)\\
&{}{-}\, 2 (\log L)_- (\psi_1)
-2c_i\sum_{j=0}^{i-1} \lambda^{i-1-j} \bigl(N_j-R_j w\Lambda^{-1}\bigr) (\psi_1).\nn
\end{align*}
Using these relations and using \eqref{Laxa1bi} and \eqref{Laxa2bi} with $i=0$, and again by a lengthy calculation, we find
\begin{align}
& {\rm e}^{-\Lambda^{-1}(\sigma)}\frac{\p \bigl({\rm e}^{\Lambda^{-1}(\sigma)}d\bigr)}{\p b_i}
= - 2\Lambda^{-1}(r_0)d + \p_x(d) = 0. \nn
\end{align}
The lemma is proved.
\end{proof}

It follows from Lemma~\ref{keylemma} and \eqref{psiphiform1} and \eqref{psiphiform2} that
\[
d(\lambda; x, {\bf b}, {\bf m};\e) = \lambda {\rm e}^{-\Lambda^{-1}(\sigma(x, {\bf b},{\bf m};\e))} \sum_{k\ge0} d_k \lambda^{-k},
\]
where $d_k$, $k\ge0$, are constants with $d_0=1$.
Therefore, for any choice $\psi_1$ of
the wave functions of type~A associated to $(v(x, {\bf b}, {\bf m};\e), w(x, {\bf b}, {\bf m};\e))$, there exists a choice, $\psi_2$, of
the wave functions of type~B associated to $(v(x, {\bf b}, {\bf m};\e), w(x, {\bf b}, {\bf m};\e))$, such that
\beq\label{dformula1002}
d(\lambda; x, {\bf b}, {\bf m};\e) = \lambda {\rm e}^{-\Lambda^{-1}(\sigma(x, {\bf b},{\bf m};\e))}.
\eeq
As a generalization of the terminology in~\cite{Y}, we say that $\psi_1$, $\psi_2$ form {\it a pair of wave functions}
associated to $(v({\bf b}+x{\bf 1}, {\bf m};\e), w({\bf b}+x{\bf 1}, {\bf m};\e))$ if~\eqref{dformula1002} holds.

We now define
\begin{align*}%\label{dpsiab}
&D(\lambda,\mu; x, {\bf b}, {\bf m};\e) \\
&\qquad{}:= \frac{\psi_1(\lambda; x, {\bf b}, {\bf m};\e) \psi_2(\mu; x-\e, {\bf b}, {\bf m};\e)
- \psi_1(\lambda; x-\e, {\bf b}, {\bf m};\e) \psi_2(\mu; x, {\bf b}, {\bf m};\e)}{\lambda-\mu}, \nn
\end{align*}
which coincides with the one introduced in~\cite{Y} when restricted to $b_2=b_3=\dots=0$.


We arrive at the following theorem.
\begin{Theorem} The following identity holds:
\begin{align}
& \e^n \sum_{i_1,\dots,i_n\ge0}
\prod_{j=1}^n \frac{(i_j+1)!}{\lambda_j^{i_j+2}} \frac{\p^n \log \tau({\bf b}+x{\bf 1},{\bf m}; \e)}{\p m_{i_1} \cdots \p m_{i_n}}
\nn\\
&\qquad{}= (-1)^{n-1}
\sum_{\text{$n$-cycles}}
\prod_{j=1}^n B\bigl(\lambda_{\sigma(j)},\lambda_{\sigma(j+1)}; x, {\bf b}, {\bf m};\e\bigr) - \frac{\delta_{n,2}}{(\lambda-\mu)^2},\label{Y926-ext-b}
\end{align}
where $n\ge2$ and
\begin{align*}%\label{eqn:hat-A-29-ext}
&B(\lambda,\mu;x, {\bf b},{\bf m};\e)\\
&\qquad{} := \frac{{\rm e}^{\sigma(x-\e, {\bf b},{\bf m};\e)}}{\mu} D(\lambda,\mu; x, {\bf b}, {\bf m};\e)
\frac
{{\rm e}^{\sum_{i\ge0} \frac2{i!} \frac{b_i}\e \mu^i (\log \mu-c_i)+\sum_{i\ge0} \frac{m_i}{(i+1)!\e} \mu^{i+1}}}
{{\rm e}^{\sum_{i\ge0} \frac2{i!} \frac{b_i}\e \lambda^{i} (\log \lambda-c_i)+\sum_{i\ge0} \frac{m_i}{(i+1)!\e} \lambda^{i+1}}}\frac{\mu^{\frac{x}\e}}{\lambda^{\frac{x}\e}}. \nn
\end{align*}
\end{Theorem}
\begin{proof}
Let
\[
\Psi_{\rm pair}(\lambda; x, {\bf b}, {\bf m};\e) :=
\begin{pmatrix} \psi_1(\lambda; x, {\bf b}, {\bf m};\e) & \psi_2(\lambda; x, {\bf b}, {\bf m};\e)
\\
\psi_1(\lambda; x-\e, {\bf b}, {\bf m};\e) & \psi_2(\lambda; x-\e, {\bf b}, {\bf m};\e)
\end{pmatrix} .
\]
By using the arguments the same as in~\cite{Y}, one can show that
\beq\label{RPsi}
R(\lambda; x, {\bf b}, {\bf m};\e) = \Psi_{\rm pair}(\lambda; x, {\bf b}, {\bf m};\e) \begin{pmatrix} 1 & 0 \\ 0 & 0\end{pmatrix}
\Psi_{\rm pair}(\lambda; x, {\bf b}, {\bf m};\e)^{-1}
\eeq
(see~\cite[the proof of Proposition~3]{Y}).
It follows from~\eqref{multiderivativesresolvent} and \eqref{RPsi} that for $n\ge2$,
\begin{align*}%\label{Y421-1-ext}
&\e^n\sum_{i_1,\dots,i_n\ge0} \prod_{j=1}^n \frac{(i_j+1)!}{\lambda_j^{i_j+2}}
\frac{\p^n \log \tau({\bf b}+x{\bf 1},{\bf m})}{\p m_{i_1} \cdots \p m_{i_n}} \\
&\qquad{}= (-1)^{n-1}
\frac{{\rm e}^{n \sigma(x-\e,{\bf b}, {\bf m};\e)}}{\prod_{j=1}^n \lambda_j} \sum_{\text{$n$-cycles}}
\prod_{j=1}^n D\bigl(\lambda_{\sigma(j)},\lambda_{\sigma(j+1)}; x, {\bf b}, {\bf m};\e\bigr) - \frac{\delta_{n,2}}{(\lambda-\mu)^2}. \nn
\end{align*}
This yields identity~\eqref{Y926-ext-b}.
\end{proof}

By taking ${\bf m}={\bf 0}$, we immediately obtain the following corollary.
\begin{Corollary} We have
\begin{align}
& \e^n\sum_{i_1,\dots,i_n\ge0}
\prod_{j=1}^n \frac{(i_j+1)!}{\lambda_j^{i_j+2}} \frac{\p^n \log \tau({\bf b}+x{\bf 1},{\bf m};\e)}{\p m_{i_1} \cdots \p m_{i_n}}\bigg|_{{\bf m}={\bf 0}}
\nn\\
& \qquad{} = (-1)^{n-1}
\sum_{\text{$n$-cycles}}
\prod_{j=1}^n B\bigl(\lambda_{\sigma(j)},\lambda_{\sigma(j+1)}; x, {\bf b};\e\bigr) - \frac{\delta_{n,2}}{(\lambda-\mu)^2},\label{Y926-ext}
\end{align}
where $B(\lambda,\mu; x, {\bf b};\e):=B(\lambda,\mu; x, {\bf b},{\bf m}={\bf 0};\e)$.
\end{Corollary}

In terms of the functions $\phi_1$ and $\phi_2$, the pair condition~\eqref{dformula1002} reads
\begin{gather}
\phi_1(\lambda; x, {\bf b}, {\bf m};\e) \phi_2(\lambda; x-\e, {\bf b}, {\bf m};\e)\nn \\
\qquad{}- w(x, {\bf b}, {\bf m};\e)
\lambda^{-2} \phi_2(\lambda; x, {\bf b}, {\bf m};\e) \phi_1(\lambda; x-\e, {\bf b}, {\bf m};\e) \equiv 1,\label{pairconditionphi}
\end{gather}
and the function $B(\lambda,\mu;{\bf b},{\bf m};\e)$ reads
\begin{align}
&B(\lambda,\mu; x, {\bf b}, {\bf m};\e) \label{Bphiab} \\
&\quad{}:= \frac{\phi_1(\lambda; x, {\bf b}, {\bf m};\e) \phi_2(\mu; x-\e,{\bf b}, {\bf m};\e)
- \frac{w(x, {\bf b}, {\bf m};\e)}{\lambda\mu} \phi_1(\lambda; x-\e,{\bf b}, {\bf m};\e) \phi_2(\mu; x, {\bf b}, {\bf m};\e)}{\lambda-\mu}. \nn
\end{align}

\begin{Lemma}
The function $B(\lambda,\mu; x, {\bf b},{\bf m};\e)$ admits the following expansion:
\beq\label{eqn:hat-A-26-extwitht}
B(\lambda,\mu; x, {\bf b},{\bf m};\e) =
\frac1{\lambda-\mu} + \sum_{i,j\ge0} \frac{A_{i,j}(x, {\bf b},{\bf m};\e)}{\lambda^{j+1} \mu^{i+1}} .
\eeq
In particular,
\beq\label{eqn:hat-A-26-ext}
B(\lambda,\mu;x, {\bf b};\e) =
\frac1{\lambda-\mu} + \sum_{i,j\ge0} \frac{A_{i,j}(x, {\bf b};\e)}{\lambda^{j+1} \mu^{i+1}} ,
\eeq
where
$A_{i,j}(x, {\bf b};\e)=A_{i,j}(x, {\bf b},{\bf m}={\bf 0};\e)$.
\end{Lemma}
\begin{proof}
Recall the following identity:
\begin{align*}
\phi_2(\mu) = \phi_2(\lambda)+\phi_2'(\lambda)(\mu-\lambda)
+ (\mu-\lambda)^2 \p_\lambda \left(\frac{\phi_2(\lambda)-\phi_2(\mu)}{\lambda-\mu}\right),
\end{align*}
where we omit the arguments $x$, ${\bf b}$, ${\bf m}$ from $\phi_2(\lambda; x, {\bf b}, {\bf m};\e)$.
Similarly,
\[
\frac{\phi_2(\mu)}{\mu} = \frac{\phi_2(\lambda)}{\lambda}+\p_\lambda\biggl(\frac{\phi_2(\lambda)}{\lambda}\biggr) (\mu-\lambda)
+ (\mu-\lambda)^2 \p_\lambda \biggl(\frac{\phi_2(\lambda)/\lambda-\phi_2(\mu)/\mu}{\lambda-\mu}\biggr).
\]
Substituting these identities in~\eqref{Bphiab} and using \eqref{pairconditionphi}, \eqref{phiexpansion1} and \eqref{phiexpansion2}, we find the validity of the expansion~\eqref{eqn:hat-A-26-extwitht}.
\end{proof}


\begin{proof}[Proof of Theorem~\ref{thm:toda-kp-ext}]
By using \eqref{Y926-ext} and \eqref{eqn:hat-A-26-ext} and Corollary~\ref{conversthmA}.
\end{proof}
